% GENERATED FILE. Edit canonical lessons, then run npm run generate:books.
% canonical-source-hash: fnv1a64:93572a47c3163cbb
% canonical-lessons: GE-C502-angeln, GE-C502-koeder, GE-C502-anbeissen, GE-C502-objektiv, GE-C502-knipsen

\chapter{Angeln und fotografieren}
\label{ch:ge-angeln-und-fotografieren}
\hlchaptermodality{502}

\begin{chapteropening}
\textbf{By the end of this chapter:} \emph{I can talk about fishing and taking photos.}

Complete the last lesson of chapter 502: I can talk about fishing and taking photos.
\end{chapteropening}

\section[angeln]{angeln — to fish (with a rod)}

\label{lesson:GE-C502-angeln}



\begin{quote}
Before the new one: say the German for to cheat, then the German for round.
\end{quote}

\subsection*{You'll want to know: angeln}
\textbf{angeln} — \textquotedblleft{}to fish (with a rod)\textquotedblright{}.

\begin{quote}
Am Samstag angelt mein Vater am See.
\end{quote}

\begin{quote}
\emph{On Saturday my father goes fishing at the lake.}
\end{quote}

\emph{Angeln gehen} — to go fishing. \emph{Die Angel} — fishing rod.

\begin{cousinweb}
\emph{angeln} is built on \emph{Angel} \textquotedblleft{}fishing rod\textquotedblright{}, inherited from Middle High German \emph{angel} and Old High German \emph{angul} \textquotedblleft{}fish hook\textquotedblright{}, from Proto-Germanic \emph{angulō} (EtymDB).
\end{cousinweb}

\begin{grammarlens}[title={What you've built}]
One more word for telling, explaining and describing at B1.
\end{grammarlens}

\subsection*{Your turn}
\begin{itemize}
\raggedright
  \item \emph{Say it:} \emph{angeln}
  \item \emph{Say it:} \emph{angeln}, once more
  \item \emph{Say it:} a sentence of your own with \emph{angeln}
  \item \emph{Recall:} say the German for to cheat, then the German for round, then say \emph{angeln} again
\end{itemize}

\begin{tcolorbox}[breakable,colback=teal!4,colframe=teal!35!black,title={Before you move on}]
What does angeln mean? (\textquotedblleft{}To fish (with a rod)\textquotedblright{}.) Say it once more.
\end{tcolorbox}
\section[der Köder]{der Köder — bait}

\label{lesson:GE-C502-koeder}



\begin{quote}
Before the new one: say the German for round, then the German for to fish.
\end{quote}

\subsection*{You'll want to know: der Köder}
\textbf{der Köder} — \textquotedblleft{}bait\textquotedblright{}.

\begin{quote}
Als Köder nimmt er Würmer.
\end{quote}

\begin{quote}
\emph{He uses worms as bait.}
\end{quote}

Plural: \emph{die Köder}.

\begin{cousinweb}
\emph{Köder} is inherited from Middle High German \emph{querder} and Old High German \emph{querdar}, from Proto-Germanic \emph{kwerþraz} \textquotedblleft{}bait\textquotedblright{} (EtymDB).
\end{cousinweb}

\begin{grammarlens}[title={What you've built}]
One more word for telling, explaining and describing at B1.
\end{grammarlens}

\subsection*{Your turn}
\begin{itemize}
\raggedright
  \item \emph{Say it:} \emph{der Köder}
  \item \emph{Say it:} \emph{der Köder}, once more
  \item \emph{Say it:} a sentence of your own with \emph{der Köder}
  \item \emph{Recall:} say the German for round, then the German for to fish, then say \emph{der Köder} again
\end{itemize}

\begin{tcolorbox}[breakable,colback=teal!4,colframe=teal!35!black,title={Before you move on}]
What does der Köder mean? (\textquotedblleft{}Bait\textquotedblright{}.) Say it once more.
\end{tcolorbox}
\section[anbeißen]{anbeißen — to bite (of a fish); to take a bite of}

\label{lesson:GE-C502-anbeissen}



\begin{quote}
Before the new one: say the German for to fish, then the German for bait.
\end{quote}

\subsection*{You'll want to know: anbeißen}
\textbf{anbeißen} — \textquotedblleft{}to bite (of a fish); to take a bite of\textquotedblright{}.

\begin{quote}
Zwei Stunden lang hat kein Fisch angebissen.
\end{quote}

\begin{quote}
\emph{For two hours not a single fish bit.}
\end{quote}

Separable: \emph{ein Fisch beißt an}. Also to bite into: \emph{einen Apfel anbeißen}.

\begin{cousinweb}
\emph{anbeißen} is built on \emph{an} \textquotedblleft{}on\textquotedblright{} and \emph{beißen} \textquotedblleft{}to bite\textquotedblright{}, inherited from Old High German \emph{bīzan}, from Proto-Germanic \emph{bītaną} (EtymDB).
\end{cousinweb}

\begin{grammarlens}[title={What you've built}]
One more word for telling, explaining and describing at B1.
\end{grammarlens}

\subsection*{Your turn}
\begin{itemize}
\raggedright
  \item \emph{Say it:} \emph{anbeißen}
  \item \emph{Say it:} \emph{anbeißen}, once more
  \item \emph{Say it:} a sentence of your own with \emph{anbeißen}
  \item \emph{Recall:} say the German for to fish, then the German for bait, then say \emph{anbeißen} again
\end{itemize}

\begin{tcolorbox}[breakable,colback=teal!4,colframe=teal!35!black,title={Before you move on}]
What does anbeißen mean? (\textquotedblleft{}To bite (of a fish); to take a bite of\textquotedblright{}.) Say it once more.
\end{tcolorbox}
\section[das Objektiv]{das Objektiv — lens (of a camera)}

\label{lesson:GE-C502-objektiv}



\begin{quote}
Before the new one: say the German for bait, then the German for to bite.
\end{quote}

\subsection*{You'll want to know: das Objektiv}
\textbf{das Objektiv} — \textquotedblleft{}lens (of a camera)\textquotedblright{}.

\begin{quote}
Für Fotos von Vögeln braucht man ein gutes Objektiv.
\end{quote}

\begin{quote}
\emph{For photos of birds you need a good lens.}
\end{quote}

Plural: \emph{die Objektive}. Not the adjective \emph{objektiv} — objective, fair.

\begin{cousinweb}
\emph{Objektiv} is a learned word built on Latin \emph{obiectus} \textquotedblleft{}set in front\textquotedblright{}, from \emph{obicere} \textquotedblleft{}to throw in front\textquotedblright{}, built on \emph{iacere} \textquotedblleft{}to throw\textquotedblright{} (Lewis \& Short): the lens faces the object.
\end{cousinweb}

\begin{grammarlens}[title={What you've built}]
One more word for telling, explaining and describing at B1.
\end{grammarlens}

\subsection*{Your turn}
\begin{itemize}
\raggedright
  \item \emph{Say it:} \emph{das Objektiv}
  \item \emph{Say it:} \emph{das Objektiv}, once more
  \item \emph{Say it:} a sentence of your own with \emph{das Objektiv}
  \item \emph{Recall:} say the German for bait, then the German for to bite, then say \emph{das Objektiv} again
\end{itemize}

\begin{tcolorbox}[breakable,colback=teal!4,colframe=teal!35!black,title={Before you move on}]
What does das Objektiv mean? (\textquotedblleft{}Lens (of a camera)\textquotedblright{}.) Say it once more.
\end{tcolorbox}
\section[knipsen]{knipsen — to snap (a photo)}

\label{lesson:GE-C502-knipsen}



\begin{quote}
Before the new one: say the German for to bite, then the German for lens.
\end{quote}

\subsection*{You'll want to know: knipsen}
\textbf{knipsen} — \textquotedblleft{}to snap (a photo)\textquotedblright{}.

\begin{quote}
Im Urlaub knipst sie jeden Tag Hunderte Fotos.
\end{quote}

\begin{quote}
\emph{On holiday she snaps hundreds of photos every day.}
\end{quote}

Informal for \emph{fotografieren}. Also to punch a ticket.

\begin{cousinweb}
\emph{knipsen} is usually explained as an imitative word for a short click; EtymDB sets it beside Dutch \emph{knippen} \textquotedblleft{}to snip\textquotedblright{}.
\end{cousinweb}

\begin{grammarlens}[title={What you've built}]
One more word for telling, explaining and describing at B1.
\end{grammarlens}

\subsection*{Your turn}
\begin{itemize}
\raggedright
  \item \emph{Say it:} \emph{knipsen}
  \item \emph{Say it:} \emph{knipsen}, once more
  \item \emph{Say it:} a sentence of your own with \emph{knipsen}
  \item \emph{Recall:} say the German for to bite, then the German for lens, then say \emph{knipsen} again
\end{itemize}

\begin{tcolorbox}[breakable,colback=teal!4,colframe=teal!35!black,title={Before you move on}]
What does knipsen mean? (\textquotedblleft{}To snap (a photo)\textquotedblright{}.) Say it once more.
\end{tcolorbox}
